% Módulo científico sin preámbulo ni portada autónoma.
% Dependencia causal: atlas, operadores K_D/K_\Omega y contraste metrológico.

\chapter{Réalisation équivariante des espèces et opérateurs de masse : fibres généalogiques, mesures spectrales, pôles complexes et compositions matérielles}
\label{hmtmasscl:chap-realizacion-equiv}

\section{Domaine interne de la réalisation matérielle et classification de ses codomaines}

\subsection{Séparation entre ligne spectrale, multiplet, distribution, pôle complexe, objet catégorique et critère d'impossibilité}

Soit \(\mathcal P\) l'ensemble des états physiques individualisés. Un état
élémentaire est représenté par le descripteur
\[
 X=(\mathfrak c,\rho,g,f,q,\epsilon,\varsigma,\mathfrak a),
\]
qui conserve la classe de représentation \(\mathfrak c\), la représentation interne
\(\rho\), la génération \(g\), la saveur \(f\), la charge \(q\), l'orientation de feuillet
\(\epsilon\), le secteur statistique \(\varsigma\) et les données de couplage
\(\mathfrak a\). Un état composé ajoute des constituants ordonnés, un arbre de
composition, une frontière commune et des nombres quantiques collectifs. Une résonance
ajoute les canaux ouverts et la prescription causale de la résolvante.

La réalisation physique doit répondre à trois questions distinctes :

\begin{enumerate}
\item quelle fibre, orbite ou multisection de l'atlas correspond au descripteur ;
\item quel opérateur positif publie la masse sur cette fibre ;
\item quel couplage produit un scalaire, un multiplet, une distribution ou
      un pôle complexe.
\end{enumerate}

Confondre ces questions introduit un sélecteur rétrospectif. La réponse
correcte s'obtient par une relation naturelle, un opérateur de fibre et une
compression spectrale, dans cet ordre.

\subsection{Antériorité de l'atlas généalogique par rapport à toute identification métrologique}

La Loi générale des masses ne part pas d'une table de noms ni d'une liste de
valeurs expérimentales. Son domaine interne est la chaîne finie
\[
 \mathcal R_{324}\longrightarrow \mathcal C_{81}
 \longrightarrow \mathcal F_{56}\longrightarrow \mathcal M_{13},
\]
où \(\mathcal R_{324}\) contient les \(4\) orientations de chacune des
\(81\) cellules, \(\mathcal F_{56}\) réunit des histoires de routes équivalentes et
\(\mathcal M_{13}\) conserve les multisections opératorielles. Le secteur nul
bifurque avant l'application du caractère positif. Par conséquent, ni une table
historique abrégée ni une taxonomie de noms physiques ne peuvent remplacer ce
domaine.

\section{Produit fibré entre descripteurs d'espèce et voies de l'atlas 81/56/13/324}

\subsection{Condition exacte de non-vacuité et sélection prospective sans sélecteur ponctuel illicite}

Soit \(\mathcal D\) l'espace fini des données internes déjà construites avant la
confrontation : représentation, orientation, mémoire, frontière, charge, saveur,
couleur, génération et secteur statistique. Notons
\[
 d_{\mathcal P}:\mathcal P\longrightarrow\mathcal D,
 \qquad
 d_{\mathcal R}:\mathcal R_{324}\longrightarrow\mathcal D
\]
les applications de descripteur physique et de descripteur généalogique. On définit l'objet
de réalisation par le produit fibré
\[
 \mathcal J
 =\mathcal P\times_{\mathcal D}\mathcal R_{324}
 =\{(X,\gamma):d_{\mathcal P}(X)=d_{\mathcal R}(\gamma)\}.
\]
La fibre d'une espèce est
\[
\mathscr S(X)=\{\gamma\in\mathcal R_{324}:(X,\gamma)\in\mathcal J\}.
\]
Un descripteur (X) est dit réalisable par l'atlas lorsque
\[
 d_{\mathcal P}(X)\in\operatorname{im}d_{\mathcal R}.
\]
Cette condition équivaut à \(\mathscr S(X)\ne\varnothing\). Elle ne se déduit pas
de la simple écriture du produit fibré : elle constitue la condition exacte d'
existence qui doit être vérifiée pour chaque descripteur physique.
La valeur de \(\mathscr S\) peut être une route, une orbite ou une multisection.
On ne force pas une section ponctuelle lorsque la symétrie interne exige de conserver
une multiplicité.

\subsection{Équivariance des fibres, orbites et multisections sous raffinement}

\begin{definition}[Relation de réalisation admissible]
La relation \(\mathscr S\) est admissible lorsqu'elle satisfait simultanément :
\begin{enumerate}
\item elle dépend uniquement de données construites avant l'ouverture de la table extérieure ;
\item elle commute avec la conjugaison particule--antiparticule ;
\item elle est équivariante sous l'inversion cellulaire et les actions internes de
      jauge ;
\item elle est compatible avec le raffinement, le retour nonadique et la lecture
      dodécaphasique ;
\item elle conserve toute multiplicité non éliminée par un couplage physique ;
\item elle reste invariante lorsqu'on permute ou remplace les valeurs extérieures de
      masse et de largeur.
\end{enumerate}
\end{definition}

\begin{theorem}[Réalisation prospective par produit fibré]
Si \(d_{\mathcal P}\) et \(d_{\mathcal R}\) sont équivariantes par rapport à une
symétrie \(G\) et ne dépendent pas de grandeurs extérieures, alors \(\mathcal J\)
est un \(G\)-sous-objet de \(\mathcal P\times\mathcal R_{324}\). Pour chaque
descripteur réalisable \(X\), la fibre non vide \(\mathscr S(X)\) est
\(G_X\)-stable et l'assignation \(X\mapsto\mathscr S(X)\) est indépendante du
résultat auquel elle sera ultérieurement confrontée.
\end{theorem}

\begin{proof}
Pour \((X,\gamma)\in\mathcal J\), on a
\(d_{\mathcal P}(X)=d_{\mathcal R}(\gamma)\). L'équivariance donne
\[
 d_{\mathcal P}(gX)=g\,d_{\mathcal P}(X)
 =g\,d_{\mathcal R}(\gamma)=d_{\mathcal R}(g\gamma),
\]
et, par conséquent, \((gX,g\gamma)\in\mathcal J\). Si \(g\in G_X\), la
seconde composante reste dans \(\mathscr S(X)\). L'indépendance des
valeurs extérieures est immédiate, car aucune d'elles n'appartient au domaine
des deux applications du produit fibré.
\end{proof}

\section{Classification spectrale intrinsèque des espèces matérielles}

\subsection{Décomposition par support ponctuel, bloc irréductible et mesure image de fibre}

Soit \(\mathcal A_{\rm rut}\) l'algèbre commutative finie engendrée par les
observables internes de route : signature \(\mathbb Z^5\), caractère, chronologie,
\(K_D\), \(K_\Omega\), orientation, mémoire et lecteurs compatibles. Dans une
fibre \(F\), on définit
\[
 \gamma\sim_{\rm sp}\eta
 \quad\Longleftrightarrow\quad
 a(\gamma)=a(\eta)\quad\text{pour tout }a\in\mathcal A_{\rm rut}.
\]
Pour chaque classe \(\lambda\in F/\!\sim_{\rm sp}\),
\[
 P_\lambda=\sum_{\gamma\in\lambda}
 |\gamma\rangle\langle\gamma|
\]
est un projecteur spectral ; les \(P_\lambda\) sont orthogonaux et
\(\sum_\lambda P_\lambda=I_F\).

\begin{theorem}[Individualisation interne maximale]
Tout classificateur construit exclusivement avec des observables de
\(\mathcal A_{\rm rut}\) se factorise de manière unique par le quotient
\(F/\!\sim_{\rm sp}\). Par conséquent, deux descripteurs de signatures complètes
distinctes appartiennent à des secteurs distinguables ; deux descripteurs qui coïncident
pour tous les générateurs de l'algèbre ne peuvent pas être séparés par une règle
construite avec ces mêmes observables.
\end{theorem}

\begin{proof}
Les caractères de l'algèbre commutative finie séparent exactement son
spectre conjoint. Un classificateur fonctionnel de ses générateurs est constant
sur chaque fibre de l'application spectrale conjointe et se factorise donc par le
quotient. Si deux routes ont le même caractère sur toute l'algèbre, aucun
élément de celle-ci ne peut les séparer.
\end{proof}

\subsection{Distinction entre masse scalaire, multiplet et distribution spectrale}

Ce théorème clôt la limite logique de la classification. Lorsqu'une
représentation physique exige de distinguer deux états appartenant à une même
classe spectrale, elle doit exhiber un observable supplémentaire ou un couplage qui
étend \(\mathcal A_{\rm rut}\). L'absence d'un tel observable n'autorise pas à
utiliser la masse extérieure comme étiquette de sélection.

\section{Raffinement naturel du registre enrichi dans les tours positives \texorpdfstring{\(90/120\)}{90/120}}

\subsection{Énumération positive du semi-groupe \texorpdfstring{\(\langle90,120\rangle\)}{<90,120>} et exclusion du mode 0}

L'existence du semi-groupe \(\langle90,120\rangle\) ne suffit pas à elle seule :
il faut construire l'application qui assigne des coefficients à une route enrichie. Soit
\[
 \mathcal H_{90,120}
 =\{90a+120b:a,b\in\mathbb N_0,\ (a,b)\ne(0,0)\}
 =\{90,120\}\cup\{30r:r\geq6\}
 =\{h_0<h_1<h_2<\cdots\}
\]
son énumération strictement croissante, où \(\mathbb N_0=\{0,1,2,\ldots\}\),
et soit \(s_{h_n}>0\) la normalisation déclarée du mode correspondant. Le
registre compatible d'une route
\(\Gamma\) fournit une suite de données entières finies
\[
 \mathcal L(\Gamma)=(\ell_0(\Gamma),\ell_1(\Gamma),\ldots),
\]
dans laquelle chaque \(\ell_n\) réunit le retour, la retenue, l'enroulement, la mémoire et la frontière
au raffinement \(n\). Pour fixer l'encodage sans ambiguïté, soit
\[
 \nu(z)=
 \begin{cases}2z,&z\geq0,\\-2z-1,&z<0,\end{cases}
 \qquad
 \pi_{\rm C}(u,v)=\frac{(u+v)(u+v+1)}2+v,
\]
et définissons \(\iota:\mathbb Z^d\to\mathbb N_0\) par application récursive de
\(\pi_{\rm C}\) à \(\nu(z_1),\ldots,\nu(z_d)\). Alors
\[
 q_n(\Gamma)
 =\frac{\iota(\ell_n(\Gamma))}
 {1+|\iota(\ell_n(\Gamma))|}\in(-1,1).
\]
La réalisation minimale de tour est
\[
 \eta_{h_n}(\Gamma)
 =\frac{3^{-(n+1)}q_n(\Gamma)}{s_{h_n}},
 \qquad
 \mathfrak T(\Gamma)
 =(\eta_h(\Gamma))_{h\in\mathcal H_{90,120}}.
\]

\subsection{Convergence absolue, naturalité par préfixes et borne effective de queue}

\begin{theorem}[Convergence, naturalité et cécité métrologique]
Pour toute route enrichie,
\[
 \sum_{h\in\mathcal H_{90,120}}|\eta_h(\Gamma)s_h|
 \leq\sum_{n\ge0}3^{-(n+1)}=\frac12.
\]
Si le raffinement prolonge le registre par préfixes, \(\mathfrak T\) est naturel
par rapport aux troncatures. De plus, aucune valeur de masse extérieure n'intervient
dans \(\mathfrak T\), et la queue après \(N\) modes satisfait
\[
 \sum_{n>N}|\eta_{h_n}s_{h_n}|
 \leq\frac{3^{-(N+1)}}2.
\]
\end{theorem}

\begin{proof}
On a \(|q_n|<1\) ; les deux bornes sont des sommes géométriques. La compatibilité
par préfixes conserve tous les coefficients déjà construits et ajoute ceux du
nouveau niveau, ce qui donne le carré de naturalité avec la troncature de
\(\mathcal E_{90,120}\). La formule ne contient que le registre, l'encodage
explicite et la normalisation modale déclarée.
\end{proof}

La construction précédente fixe une réalisation minimale explicite et convergente,
relative à l'encodage \(\iota\) et aux normalisations \(s_h\) déclarées.
Une dynamique qui modifie ses poids doit déclarer l'opérateur effectuant cette
modification et démontrer de nouveau la naturalité et la convergence ; elle ne peut pas
reconstruire les poids à partir d'un résidu expérimental.

\section{Opérateurs bidirectionnels \texorpdfstring{\(K_D\) et \(K_\Omega\)}{K_D et K_Omega} et mesure spectrale conjointe}

\subsection{PVM conjointe, compression en POVM et distribution scalaire induite par un état}

Pour \(F_X=\mathscr S(X)\), soit
\[
 \mathcal H_X=\ell^2(F_X)\otimes V_{\rho_X}
\]
l'espace de réalisation, où \(V_{\rho_X}\) porte la représentation
interne correspondante. Chaque route \(\gamma\in F_X\) possède une signature
\[
 \sigma(\gamma)
 =(n_A,s_C,k_{\Delta_4},b_{120},b_\zeta)\in\mathbb Z^5
\]
et des coefficients de raffinement
\(\eta(\gamma)\in\mathcal E_{90,120}\). L'opérateur de base est
\[
 \widehat M_X^{(0)}
 =\sum_{\gamma\in F_X}
 m_{\rm clk}\,
 \chi_{\rm full}(\sigma(\gamma),\eta(\gamma))
 |\gamma\rangle\langle\gamma|\otimes I_{V_{\rho_X}}.
\]
L'échelle \(m_{\rm clk}\) descend de la branche absolue
action--horloge--masse, y compris la décennie \(10^{-34}\) ; elle n'est pas extraite d'une
masse expérimentale.

Les lecteurs \(K_D\) et \(K_\Omega\) induisent des opérateurs diagonaux exacts
\(B_X^D\) et \(B_X^\Omega\). Les deux orientations publient
\[
 \widehat M_X^{\beta,r}
 =R_{\rm act}^{B_X^r/2}\widehat M_X^{(0)}
  R_{\rm act}^{B_X^r/2},
 \qquad r\in\{D,\Omega\}.
\]
Les deux sont positives et commutent dans la base des routes. Par le théorème spectral
conjoint, il existe une unique mesure spectrale projective
\(E_X^{D,\Omega}\) sur \(\mathbb R_+^2\) telle que
\[
 \widehat M_X^{\beta,D}
 =\int x\,dE_X^{D,\Omega}(x,y),
 \qquad
 \widehat M_X^{\beta,\Omega}
 =\int y\,dE_X^{D,\Omega}(x,y).
\]
Dans la base des routes,
\[
 E_X^{D,\Omega}(A)|\gamma\rangle
 =\mathbf 1_A\!\left(m_D(\gamma),m_\Omega(\gamma)\right)
  |\gamma\rangle.
\]
C'est la sortie bidirectionnelle canonique : elle conserve la paire complète et ne
suppose pas d'avance une fonctionnelle scalaire.

Pour un canal \(P_{X,a}\), la mesure spectrale comprimée
\[
 F_{X,a}^{D,\Omega}(A)
 =\left.P_{X,a}E_X^{D,\Omega}(A)P_{X,a}\right|_{P_{X,a}\mathcal H_X}
\]
est une POVM sur le sous-espace préparé. Pour une densité \(\rho_X\)
portée par ce sous-espace, la distribution scalaire conjointe est
\[
 \mu_{X,a}^{D,\Omega}(A)
 =\operatorname{Tr}\!\left(
 \rho_X F_{X,a}^{D,\Omega}(A)
 \right).
\]
Toute fonctionnelle déclarée \(f:\mathbb R_+^2\to\mathbb R_+\) induit
\(\mu_{X,a}^f=f_\#\mu_{X,a}^{D,\Omega}\) sans modifier la PVM d'origine.

\subsection{Moments mixtes, support spectral ponctuel et conservation des multiplicités}

Soit
\[
 \boldsymbol\beta_X(\gamma)
 =(m_D(\gamma),m_\Omega(\gamma)).
\]
La mesure image uniforme de la fibre est
\[
 \nu_X=(\boldsymbol\beta_X)_\#
 \left(\frac1{|F_X|}\sum_{\gamma\in F_X}\delta_\gamma\right)
 =\frac1{|F_X|}\sum_{\gamma\in F_X}
 \delta_{(m_D(\gamma),m_\Omega(\gamma))}.
\]
Sa transformée et ses moments mixtes sont
\[
 Z_X(s,t)=\int e^{sx+ty}\,d\nu_X(x,y),
 \qquad
 \mathfrak m_{p,q}(X)=\int x^py^q\,d\nu_X(x,y).
\]
Comme \(\nu_X\) est à support fini, \(Z_X\), ou de manière équivalente la famille
complète des moments mixtes, détermine le multiensemble conjoint des paires de
lecteurs à permutation des routes près. La mesure image est donc un
invariant complet de la paire opératorielle ; ce n'est pas une moyenne de masses.

Si \(u:F_X\to F_Y\) est une bijection qui conserve le registre, l'orientation et
les deux lecteurs, l'unitaire \(U_u|\gamma\rangle=|u\gamma\rangle\) satisfait
\[
 U_uE_X^{D,\Omega}(A)=E_Y^{D,\Omega}(A)U_u,
 \qquad
 \nu_X=\nu_Y.
\]
La même égalité vaut pour les distributions physiques lorsque la densité et le
projecteur sont transportés par \(U_u\). Cette naturalité n'est affirmée que pour
les entrelaceurs effectivement construits ; une coïncidence de
cardinaux ne crée pas à elle seule une équivalence de fibres.

\section{Scalarisation par symétrie et décomposition isotypique des fibres}

\subsection{Espérance conditionnelle de symétrie et critère de réduction à une ligne spectrale}

Soit \(G_X\) le groupe compact de symétrie non brisée de la fibre et
\(U_X:G_X\to\mathcal U(\mathcal H_X)\) sa représentation. L'espérance
conditionnelle de symétrie est définie par
\[
 \mathbb E_X(T)
 =\int_{G_X}U_X(g)TU_X(g)^*\,d\mu_X(g),
\]
où \(\mu_X\) est la mesure de Haar normalisée. L'opérateur physique de masse
avant l'ouverture des canaux est défini, lorsque le type d'état déclare une
fonctionnelle \(\Phi_X:\mathbb R_+^2\to\mathbb R_+\), par
\[
 \mathcal M_X
 =\mathbb E_X\!\left(
   \Phi_X(\widehat M_X^{\beta,D},\widehat M_X^{\beta,\Omega})
  \right).
\]
Le calcul fonctionnel s'effectue avec la mesure conjointe
\(E_X^{D,\Omega}\). Dans une représentation contenant un entrelaceur d'
échange des lecteurs, on peut prendre \(\Phi_X(x,y)=\sqrt{xy}\). Si le type
physique ne déclare pas \(\Phi_X\), la sortie complète reste la PVM conjointe ;
on n'invente pas de scalaire.

\begin{theorem}[Clôture scalaire, multiplet et distribution]
L'opérateur \(\mathcal M_X\) appartient au commutant
\(U_X(G_X)'\). Si
\[
 \mathcal H_X\simeq
 \bigoplus_\lambda V_\lambda\otimes\mathbb C^{m_\lambda},
\]
alors
\[
 \mathcal M_X=
 \bigoplus_\lambda I_{V_\lambda}\otimes M_{X,\lambda}.
\]
Dans un secteur irréductible de multiplicité \(1\), la restriction est un scalaire.
Dans un secteur réductible, la sortie exacte est le spectre des blocs
\(M_{X,\lambda}\), avec leurs multiplicités. Il n'existe pas d'indétermination :
il existe un codomaine distinct.
\end{theorem}

\begin{proof}
L'invariance de la mesure de Haar implique
\(U_X(h)\mathbb E_X(T)U_X(h)^*=\mathbb E_X(T)\) pour tout \(h\in G_X\).
La décomposition découle du théorème de complète réductibilité et du lemme de
Schur. Le cas irréductible produit un multiple de l'identité ; le cas
réductible conserve exactement les blocs de multiplicité.
\end{proof}

Soit \(P_{X,a}\) le projecteur déterminé par l'état préparé et le canal ou l'
appareil physique \(a\). La mesure spectrale observable est
\[
 \mu_{X,a}(B)=
 \langle\Psi_X,
 P_{X,a}E_{\mathcal M_X}(B)P_{X,a}\Psi_X\rangle.
\]
Il y a une ligne scalaire précisément lorsque le support spectral comprimé est
ponctuel :
\[
 E_{\mathcal M_X}(\{m_X\})P_{X,a}=P_{X,a},
 \qquad\text{de manière équivalente}\qquad
 (\mathcal M_X-m_X I)\,P_{X,a}=0.
\]
L'égalité isolée
\(P_{X,a}\mathcal M_X P_{X,a}=m_X P_{X,a}\) ne fixe qu'une espérance si le
sous-espace n'est pas réducteur et ne suffit pas à démontrer une ligne spectrale.
Si la compression conserve plusieurs valeurs propres, la sortie est un multiplet ou une
distribution. La remplacer par une moyenne numérique détruirait l'information
physique et généalogique.

\subsection{Moyenne géométrique des lecteurs uniquement sous invariance d'échange démontrée}

Lorsque la représentation physique exige une symétrie sous l'échange
\(D\leftrightarrow\Omega\), l'agrégat opératoriel minimal est la moyenne
géométrique
\[
 \widehat M_X^{\rm sym}
 =\widehat M_X^{\beta,D}\#\widehat M_X^{\beta,\Omega},
 \qquad
 A\#B
 =A^{1/2}\bigl(A^{-1/2}BA^{-1/2}\bigr)^{1/2}A^{1/2},
\]
sur le support positif commun. Dans le cas diagonal,
\[
 \widehat M_X^{\rm sym}
 =R_{\rm act}^{(B_X^D+B_X^\Omega)/4}
  \widehat M_X^{(0)}
  R_{\rm act}^{(B_X^D+B_X^\Omega)/4}.
\]
Cette formule est une conséquence de l'hypothèse explicite d'échange ; elle ne
remplace pas la mesure conjointe lorsque cette symétrie n'a pas été établie.

\begin{theorem}[Caractérisation de l'agrégat symétrique]
La moyenne géométrique est positive, symétrique, covariante par congruence pour
toute \(S\) inversible et autoduale :
\[
 A\#B=B\#A,
 \qquad
 S^*(A\#B)S=(S^*AS)\#(S^*BS),
 \qquad
 (A\#B)^{-1}=A^{-1}\#B^{-1}.
\]
De plus, \(A\#A=A\). Par conséquent, dans les secteurs où il existe un
entrelaceur physique qui échange les lecteurs, l'agrégat n'introduit pas de
paramètre et ne privilégie aucune orientation.
\end{theorem}

\begin{proof}
L'identité de congruence s'entend pour \(S\) inversible. Les identités
résultent du calcul fonctionnel de la racine positive de
\(A^{-1/2}BA^{-1/2}\). Dans le cas commutatif, elles se réduisent à
\(A\#B=(AB)^{1/2}\).
\end{proof}

Il existe en outre un agrégat scalaire canonique sous des axiomes distincts. Si
\(s_n\) est continu, invariant par permutation, multiplicatif coordonnée par
coordonnée et normalisé par \(s_n(cI)=c\), alors
\[
 s_n(T)=\det(T)^{1/n}.
\]
En effet, en passant aux logarithmes, la continuité et l'additivité rendent la
fonctionnelle linéaire sur les logarithmes des valeurs propres ; la symétrie égalise tous
les coefficients et la normalisation fixe leur valeur à \(1/n\). Pour la paire de
lecteurs, la symétrie d'échange donne l'invariant
\[
 S_{D,\Omega}
 =\sqrt{
 \operatorname{ndet}(T_D)\operatorname{ndet}(T_\Omega)},
 \qquad
 \operatorname{ndet}(T)=\det(T)^{1/n}.
\]
Ce résultat classifie la fibre complète sous les axiomes énoncés. Il ne s'
identifie pas automatiquement à une ligne de masse : cette identification exige
que le couplage physique adopte précisément cette fonctionnelle.

\section{Secteur nul et antériorité de la carte de masse nulle par rapport au caractère positif}

\subsection{Séparation des représentations photonique \texorpdfstring{\(U(1)\)}{U(1)} et gluonique adjointe \texorpdfstring{\(SU(3)\)}{SU(3)}}

Le photon appartient à la carte nulle électromagnétique et les gluons à la carte
nulle adjointe de couleur. Leurs différentes représentations, polarisations et
contraintes de jauge restent visibles bien que les deux masses soient nulles.

\subsection{Non-accessibilité du secteur nul comme limite de la branche positive de masse}

La nullité ne s'obtient pas en prenant une limite du caractère positif. Le domaine
se sépare d'abord en
\[
 \mathcal H_{\rm phys}=\mathcal H_0\oplus\mathcal H_+,
\]
avec
\[
 \mathcal M|_{\mathcal H_0}=0,
 \qquad
 \mathcal M|_{\mathcal H_+}>0.
\]

\section{Clôture bêta de l'électron et secteurs leptoniques chargés}

\subsection{Unicité de l'électron dans la cellule \texorpdfstring{\((5,5)\)}{(5,5)} et préservation de la masse bêta directrice}

L'électron constitue l'exception centrale. L'inversion cellulaire possède un
unique point fixe, la cellule \((5,5)\), et sa fibre physique est donc de rang
un. Les autres multisections stables n'admettent pas, en général, de sélecteur
ponctuel équivariant. Cette différence n'est pas une lacune de l'atlas : c'est la
distinction mathématique entre une section centrale et une espèce réalisée comme
orbite ou multisection.

Pour la fibre électronique de rang un, la mesure conjointe est ponctuelle, les deux
lecteurs partagent la première composante \(80\), la compression est scalaire et l'on
retrouve la clôture bêta
\[
 m_e^\beta
 =0.5109989506900008086082571648\ldots\ {\rm MeV}.
\]
La valeur de base demeure une étape de construction ; elle ne remplace pas la clôture
bêta.

\subsection{Réalisations du muon et du tau au moyen de secteurs chargés et d'un critère spectral de scalarité}

Dans le secteur leptonique chargé, les descripteurs distinguent la fibre centrale,
la multisection de profondeur \(G_2\) et la multisection de profondeur
\(G_4\). La première produit l'unique ligne électronique ; les deux autres
produisent les opérateurs associés aux réalisations muonique et tauonique. Si
le couplage préserve plus d'un bloc, la prédiction correcte est son
spectre ; l'étiquette physique n'autorise pas à choisir l'une de ses valeurs propres par
proximité numérique.

\section{Algèbre chromatique \texorpdfstring{\(SU(3)\)}{SU(3)}, saveurs quark et transport de la masse courante}

\subsection{Algèbre qutrit de Weyl, réduction de Reynolds et projecteurs singulets}

L'orbite chromatique se réalise dans le qutrit de couleur. Les opérateurs de Weyl
engendrent \(\mathfrak{sl}_3(\mathbb C)\) ; la sélection antihermitienne de trace
nulle produit \(\mathfrak{su}(3)\), et la connexion finie transporte les trois
cartes chromatiques. La masse de saveur doit être invariante sous cette action.

Soit \(U_{\rm col}:SU(3)\to\mathcal U(V_{\rm col})\). La réduction chromatique
est l'espérance de Reynolds
\[
 \mathbb E_{\rm col}(T)
 =\int_{SU(3)}U_{\rm col}(g)T U_{\rm col}(g)^*\,dg.
\]
Dans la représentation fondamentale irréductible,
\[
 \mathbb E_{\rm col}(\mathcal M_q)
 =m_q(\mu,\mathscr R)I_3.
\]
La masse publiée est donc indépendante du rouge, du vert et du bleu. Les
coordonnées \(\mu\) et \(\mathscr R\) identifient l'échelle et le schéma de
renormalisation ; ce ne sont pas des ambiguïtés cachées et elles doivent toujours accompagner le
nombre.

Pour les mésons et les baryons, les projecteurs de singulet sont
\[
 P_{q\bar q}^{(1)}
 =\frac1{3}|\Omega\rangle\langle\Omega|,
 \qquad
 |\Omega\rangle=\sum_{a=1}^{3}|a\bar a\rangle,
\]
et
\[
 P_{qqq}^{(1)}
 =|\varepsilon\rangle\langle\varepsilon|,
 \qquad
 |\varepsilon\rangle
 =\frac1{\sqrt6}\sum_{a,b,c}\varepsilon_{abc}|abc\rangle.
\]
Les deux éliminent l'étiquette chromatique observable sans effacer l'orbite interne
qui est intervenue dans la construction.

\subsection{Connexion de renormalisation et transport ordonné entre fibres d'échelle}

La dépendance en échelle est une section d'un fibré sur
\(t=\log\mu\). Si \(\gamma_m(g(t))\) est l'endomorphisme induit par la
connexion chromatique, le transport est défini par
\[
 \nabla_{\rm RG}m_q
 =\left(\frac{d}{dt}+\gamma_m(g(t))\right)m_q=0,
\]
\[
 m_q(t_2)
 =\mathcal P\exp\!\left(-\int_{t_1}^{t_2}
 \gamma_m(g(t))\,dt\right)m_q(t_1).
\]
Deux valeurs appartenant à des échelles ou à des schémas distincts ne sont pas des nombres
comparables tant qu'elles ne sont pas transportées vers la même fibre. Cette équation clôt le
type mathématique de la masse courante : ce n'est pas un scalaire absolu isolé, mais
une section transportée de manière covariante.

\section{Mélange CKM et covariance des repères spectraux quark}

\subsection{Transport unitaire de l'opérateur \texorpdfstring{\(B_d\)}{B_d} dans le repère spectral de type haut au moyen de \texorpdfstring{\(V_{\rm CKM}\)}{V_CKM}}

Le mélange CKM appartient à une application distincte. Si \(B_u\) et \(B_d\) sont les
opérateurs de masse dans les bases de type haut et de type bas, la matrice unitaire
interne satisfait
\[
 B_d^{(u)}=V_{\rm CKM} B_d V_{\rm CKM}^*.
\]

\subsection{Séparation entre la fermeture interne et la reconnaissance statistique extérieure}

Les angles dérivés de \(A\), \(C^*\) et \(\Delta_4\), la phase
\(\delta_{\rm CP}\), l'invariant de Jarlskog et l'identité du déterminant
du commutateur déterminent le mélange sans créer les valeurs propres de masse. Ainsi
restent séparés le spectre, la couleur, la génération et le changement de base.

\section{Opérateur neutrinique, mélange PMNS et critères Dirac--Majorana}

\subsection{Spectre neutre tridimensionnel et héritage de l'échelle action--horloge--masse}

Les neutrinos sont des sections fermioniques électriquement neutres. Le
caractère fermionique descend de l'holonomie spinorielle, de la complétion extérieure et
des relations canoniques d'anticommutation ; la neutralité n'altère pas leur
statistique. Le support interne est
\[
 \mathcal H_\nu=
 \bigoplus_{j=1}^{4}\mathcal H_{\nu,G_j},
 \qquad
 \dim\mathcal H_\nu=2+4+6+8=20.
\]
Les quatre profondeurs structurelles ne sont ni des saveurs ni des valeurs propres. Le
sous-espace physique tridimensionnel s'obtient au moyen du projecteur neutre
\(P_\nu\), et l'opérateur de masse est
\[
 M_\nu=P_\nu\mathcal M_\nu P_\nu.
\]
Ses trois valeurs propres ordonnées sont les masses ; l'échelle absolue provient de la
même horloge massique qui fixe les autres branches, et non des angles de mélange.

\subsection{Rang 3 de la matrice de Gram, repère polaire et critère antiunitaire de Majorana}

Soit \(M_\ell=P_\ell\mathcal M_\ell P_\ell\) l'opérateur chargé. Notons
\(F_\ell\) et \(F_\nu\) des repères spectraux orthonormaux complets, ordonnés par
les projecteurs spectraux de \(M_\ell\) et \(M_\nu\). La matrice de mélange est
\[
 U_{\rm PMNS}=F_\ell^*F_\nu.
\]

\begin{theorem}[Clôture spectrale du mélange leptonique]
Si les deux repères sont complets, \(U_{\rm PMNS}\) est unitaire. Si les
spectres sont simples, elle est déterminée à des réajustements diagonaux de phase et à des
permutations fixées par l'ordre spectral près. S'il existe une dégénérescence, le
résultat exact est une classe double
\[
 \prod_\lambda U(m_{\ell,\lambda})
 \backslash U(3) /
 \prod_\mu U(m_{\nu,\mu}),
\]
et non une matrice numérique artificiellement sélectionnée.
\end{theorem}

\begin{proof}
L'orthonormalité donne
\(U_{\rm PMNS}^*U_{\rm PMNS}=F_\nu^*F_\ell F_\ell^*F_\nu=I_3\).
Les changements de base à l'intérieur de chaque sous-espace dégénéré agissent par la
gauche ou par la droite, ce qui produit exactement la classe double
indiquée.
\end{proof}

Le registre complet détermine les éléments de \(M_\nu\) au moyen de la retenue,
du feuillet, de l'orientation, de la lacune et de l'holonomie. De manière équivalente, on peut employer
le noyau hermitien
\[
 K_{\rm reg}(c,d)
 =\sum_{r\in\{\rm car,hoj,ori,vac,hol\}}
 \frac{\langle\Xi_r(c),\Xi_r(d)\rangle}
      {\sum_x\|\Xi_r(x)\|^2},
\]
en omettant tout terme de norme nulle. Sa matrice de Gram est positive et
construit l'opérateur avant la diagonalisation. La phase de rang un
\[
 I+(e^{i\pi/6}-1)|b_K\rangle\langle b_K|
\]
demeure un contrôle négatif : son unitarité formelle ne suffit pas et son angle
\(\theta_{13}\) réfute son identification au mélange physique de rang
complet.

Soit \(G_0\) la matrice tridimensionnelle obtenue en agrégeant le noyau précédent
sur les trois secteurs neutriniques. La condition finie
\[
 \det G_0\ne0
\]
est le certificat exact que les observables incluses dans le registre
séparent trois directions. Dans ce cas, le facteur polaire
\[
 U_{G_0}=G_0(G_0^*G_0)^{-1/2}
\]
est unitaire et de rang \(3\). Si \(\det G_0=0\), le noyau déclaré ne
détermine pas de transport tridimensionnel : il faut ajouter une observable interne
indépendante et recalculer la matrice de Gram. Il n'est pas licite d'imposer
l'inversibilité en ajoutant un terme construit à partir de la matrice de mélange
que l'on cherche à obtenir. La matrice physique est fixée par les repères
spectraux de \(M_\ell\) et \(M_\nu\) ; le critère de Gram certifie la suffisance
du registre, sans sélectionner à lui seul ces repères.

Soit \(\mathcal C_\nu\) l'autoconjugaison antiunitaire du secteur neutre. Le
critère de Majorana est
\[
 \mathcal C_\nu^2=I,
 \qquad
 \mathcal C_\nu M_\nu=M_\nu\mathcal C_\nu,
 \qquad
 \mathcal C_\nu H_\nu=H_\nu\mathcal C_\nu,
\]
avec l'existence d'une base réelle fixe dans la réalification du
sous-espace physique. Lorsque ces conditions sont remplies, chaque mode peut
être représenté par \(\psi=\mathcal C_\nu\psi\). La branche de Dirac exige, en
revanche, une charge conservée \(Q\) telle que
\[
 \mathcal C_\nu Q\mathcal C_\nu^{-1}=-Q,
 \qquad Q\psi\ne0,
\]
et conserve comme modes distincts \(\psi\) et \(\mathcal C_\nu\psi\). Si aucun
des deux ensembles de conditions n'est établi, le type reste indécidé.
La neutralité, l'involution de mot ou le secteur d'Ising, pris
séparément, ne tranchent pas cette alternative.

\section{Secteur neutrinique stérile et séparation entre existence et mélange}

\subsection{Projecteur orthogonal au secteur actif et annulation du courant faible}

Une profondeur structurelle supplémentaire ne constitue pas à elle seule un neutrino
stérile. Un secteur stérile candidat est défini par un projecteur non nul
\(P_s\) qui satisfait
\[
 P_s P_{\rm act}=0,
 \qquad P_s J_{\text{faible}}=0,
\]
avec une dynamique définie sur \(P_s\mathcal H_\nu\).

\subsection{Distinction entre réduction invariante et couplage non nul avec le secteur actif}

Si de plus \([P_s,M_\nu]=0\), le secteur est réducteur et ne se mélange pas au secteur actif. Si \(P_sM_\nu P_{\rm act}\ne0\), il existe un mélange actif--stérile et nécessairement \([P_s,M_\nu]\ne0\). Par conséquent, l'existence du secteur et l'existence du mélange sont des propriétés distinctes. Si le projecteur faiblement neutre n'existe pas, le canal grossier supplémentaire est une profondeur interne du support et non une quatrième espèce.

\section{Bosons \texorpdfstring{\(W\) et \(Z\)}{W et Z} et secteur scalaire : projecteurs, résolvantes et critère de pôle}

\subsection{Opérateur effectif sectoriel et prolongement méromorphe de la résolvante réduite}

Les secteurs \(W\), \(Z\) et scalaire possèdent leurs propres descripteurs de
représentation, de route et de signature. Leur réalisation n'exige pas de les faire entrer de force dans une
fibre élémentaire de rang un. On construit d'abord l'opérateur positif de
fibre puis, dans un sous-espace \(P_X\mathcal H_X\) de dimension finie,
l'opérateur effectif de canaux
\[
 H_{X,\rm eff}(z)
 =P_X H_X P_X
 +P_X H_X Q_X(z-Q_X H_X Q_X)^{-1}Q_X H_X P_X.
\]
\subsection{Zéro isolé admissible, résidu physique et distinction par rapport à la simple ouverture de canal}

Si la résolvante réduite admet un prolongement méromorphe à travers la coupure du
continu, les pôles sont les zéros isolés, comptés avec multiplicité, de
\[
 \det\bigl(z-H_{X,\rm eff}(z)\bigr)=0,
 \qquad
 z_X=M_X c_{\rm int}^2-\frac{i}{2}\Gamma_X.
\]
Un canal ouvert ne garantit pas à lui seul l'existence de tels zéros. Lorsqu'il
existe un zéro avec le résidu physique prescrit, la partie réelle publie la masse de
pôle et la partie imaginaire la largeur. Pour le
secteur scalaire, l'identification d'une excitation exige en outre le projecteur
scalaire et ses couplages de jauge et fermioniques. Deux pôles indépendants
déterminent deux états ; une seconde étiquette sans projecteur ni pôle ne crée pas une
particule.

\section{États composés et loi cocyclique de liaison}

\subsection{Projecteurs singulets mésoniques et baryoniques et raffinement de la signature composée}

Pour des constituants \(X_1,\ldots,X_n\), l'entrée est le registre composé
\[
 \mathfrak C_{12}^{(n)}:
 (\Gamma_1,\ldots,\Gamma_n;\mathcal T,\partial)
 \longmapsto\mathcal L_{\rm comp},
\]
et non la somme des masses isolées. La composition binaire
\[
 L_1\star L_2=L_1+L_2+\mathfrak b(L_1,L_2)
\]
est associative exactement lorsque
\[
 \mathfrak b(L_1,L_2)+\mathfrak b(L_1\star L_2,L_3)
 =\mathfrak b(L_2,L_3)+\mathfrak b(L_1,L_2\star L_3).
\]
La signature, le raffinement et l'opérateur composé sont calculés après la
liaison :
\[
 \mathcal L_{\rm comp}\longmapsto\sigma_{\rm comp}
 \longmapsto\eta_{\rm comp}\longmapsto\mathcal M_{\rm comp}.
\]
Le projecteur physique réunit le singulet de couleur, le spin et la parité, l'isospin, la charge, la symétrie d'échange et le canal.

\subsection{Réduction de Feshbach--Schur et construction de l'opérateur effectif}

La masse stable est une valeur propre de la compression ; lorsque les hypothèses de prolongement méromorphe sont satisfaites, la masse et la largeur d'une résonance proviennent d'un pôle de Feshbach. Ainsi, mésons, baryons, états exotiques, noyaux, atomes et molécules possèdent le même schéma causal sans se réduire à une somme de masses.

Deux arbres de constituants représentent la même espèce s'il existe un
unitaire qui entrelace simultanément leurs opérateurs, leurs projecteurs de canal et leurs
pôles. En l'absence d'un tel unitaire, ils sont conservés comme des réalisations distinctes,
bien qu'ils partagent des constituants et des nombres quantiques partiels.

\section{Résonances, largeurs et survie des sections matérielles}

\subsection{Pôles complexes, résidus admissibles et largeurs associées au prolongement méromorphe}

Soit \(S_X\) le projecteur sur le sous-espace survivant et
\(T_X=S_X U S_X\) l'opérateur comprimé d'un pas. S'il possède une valeur propre
dominante simple
\[
 \lambda_X=r_X e^{-i\theta_X},\qquad 0<r_X<1,
\]
fixons la branche \(\theta_X\in I_\theta\), où \(I_\theta\) est l'intervalle
fondamental de quasi-énergie déterminé par la carte temporelle et son orientation.
Alors
\[
 M_X=\frac{\hbar_{\rm int}}{t_0c_{\rm int}^2}\theta_X,
 \qquad
 \Gamma_X=-\frac{2\hbar_{\rm int}}{t_0}\log r_X,
 \qquad
 \tau_X=\frac{\hbar_{\rm int}}{\Gamma_X}.
\]
\subsection{Compatibilité du raffinement généalogique avec la persistance résonante}

Sans le choix de \(I_\theta\), la phase ne détermine \(\theta_X\) que modulo
\(2\pi\) et la masse de quasi-énergie n'est pas un scalaire bien défini.
La largeur n'est pas une sixième coordonnée de la signature massique : elle provient du module
du mode ouvert. L'égalité de masse et de largeur entre particule et
antiparticule exige de transporter également les canaux et la prescription causale, et pas seulement
la parité de la signature.

\section{Complétion multiparticulaire et statistique de Fock}

\subsection{Foncteurs de Fock symétrique et extérieur induits par la parité d'holonomie}

Soit \(\mathsf{Hol}_{\pm}^{\leq1}\) la catégorie monoïdale de fibres à une
particule, avec des morphismes de raffinement contractifs, en particulier unitaires,
et une parité d'holonomie
\(\epsilon\in\{+1,-1\}\). La complétion de Fock est définie par le foncteur
\[
 \mathfrak F(H,\epsilon)=
 \begin{cases}
 \displaystyle\bigoplus_{n\ge0}\operatorname{Sym}^n(H),
     &\epsilon=+1,\\[2mm]
 \displaystyle\bigoplus_{n\ge0}\bigwedge^n\!H,
     &\epsilon=-1,
 \end{cases}
\]
et, pour un entrelaceur contractif \(T:H\to H'\),
\[
 \mathfrak F(T)=
 \bigoplus_{n\ge0}\operatorname{Sym}^n(T)
 \quad\text{ou}\quad
 \bigoplus_{n\ge0}\bigwedge^n\!T.
\]

\subsection{Relations CCR/CAR, exclusion de Pauli et domaine des morphismes non contractants}

\begin{theorem}[Transport interne de la statistique]
La complétion précédente est un foncteur d'opérateurs bornés sur
\(\mathsf{Hol}_{\pm}^{\leq1}\), compatible avec le raffinement et
monoïdal gradué. Dans la branche \(\epsilon=-1\), les opérateurs de création et
d'annihilation satisfont CAR et les opérateurs d'occupation vérifient
\(N_s^2=N_s\) ; dans la branche \(\epsilon=+1\), ils satisfont CCR. Par conséquent,
l'exclusion de Pauli et les statistiques fermionique et bosonique découlent de la
parité d'holonomie une fois fixée la représentation à une particule.
\end{theorem}

\begin{proof}
Les puissances symétriques et extérieures préservent les compositions et les identités ;
la contractivité de \(T\) garantit que leur somme directe définit un opérateur
borné. Pour les morphismes non contractifs, la seconde quantification exige un
domaine dense explicite et n'appartient pas à cette catégorie d'opérateurs bornés.
La règle des signes du produit extérieur produit CAR et l'idempotence de
l'occupation ; le quotient symétrique produit CCR. La seconde quantification des
morphismes de raffinement entrelace les opérateurs correspondants.
\end{proof}

Cette démonstration clôt le transport algébrique interne. L'identification
au théorème relativiste de spin--statistique exige, en outre, le pont
lorentzien et la localité déclarés en mécanique dimensionnelle ; les deux énoncés
ne doivent pas être confondus.

\section{Réalisation thermodynamique au moyen d’états de Gibbs et limite de pression}

\subsection{Hamiltonien fini, partition, état de Gibbs et limite cofinale de la pression}

Pour un Hamiltonien comprimé \(H_{X,V}\) sur un volume ou un raffinement
fini \(V\), l’état thermique est complètement défini par
\[
 Z_{X,V}(\beta)=\operatorname{Tr}e^{-\beta H_{X,V}},
 \qquad
 \rho_{X,V,\beta}=Z_{X,V}(\beta)^{-1}e^{-\beta H_{X,V}}.
\]
La limite thermodynamique existe lorsque la pression
\[
 p_X(\beta)=\lim_{V\nearrow\infty}
 \frac1{\beta |V|}\log Z_{X,V}(\beta)
\]
existe et est indépendante de la suite cofinale.

\subsection{Critère de condensation par saturation de la densité excitée}

La condensation bosonique est
décidée par la saturation de la densité des états excités ; elle n’est pas proclamée à
partir de la seule existence d’une complétion symétrique. Ainsi, le problème
est exprimé au moyen d’une condition mathématique vérifiable pour le
Hamiltonien et la limite, au lieu d’une attribution nominale.

\section{Secteurs topologiques, excitations virtuelles et classification de leurs codomaines}

\subsection{Catégorie tressée, données \texorpdfstring{\(N\)--\(F\)--\(R\)}{N--F--R} et équations du pentagone et de l'hexagone}

Dans un secteur topologique, le codomaine primaire est une catégorie tressée, et non
une droite de masses. Pour des objets simples \(a,b,c\), les données physiques comprennent
\[
 N_{ab}^{c},\qquad F^{abc}_{d},\qquad R^{ab}_{c},
\]
avec des équations pentagonale et hexagonale. Une masse scalaire n'existe que
lorsqu'une réalisation hamiltonienne fournit un projecteur spectral et un
gap. Les secteurs Fibonacci, Ising/Majorana et \(\mathbb Z_6\) sont ainsi
clos comme codomaines catégoriques ; aucune valeur numérique universelle ne leur est attribuée
par analogie avec une particule élémentaire.

Dans le secteur \(\mathbb Z_6\), la fusion et le tressage sont
\[
 a\otimes b=a+b\pmod 6,
 \qquad
 R_{ab}=\zeta_6^{ab},
 \qquad \zeta_6=e^{2\pi i/6}.
\]
Pour Fibonacci,
\[
 \tau\otimes\tau=\mathbf 1\oplus\tau,
 \qquad
 F_\tau^{\tau\tau\tau}=
 \begin{pmatrix}
 \varphi^{-1}&\varphi^{-1/2}\\
 \varphi^{-1/2}&-\varphi^{-1}
 \end{pmatrix},
\]
\[
 R_{\mathbf1}^{\tau\tau}=e^{-4\pi i/5},
 \qquad
 R_\tau^{\tau\tau}=e^{3\pi i/5}.
\]
Le secteur d'Ising possède des objets \(\mathbf1,\psi,\sigma\) et les règles
\[
 \psi\otimes\psi=\mathbf1,
 \qquad
 \psi\otimes\sigma=\sigma,
 \qquad
 \sigma\otimes\sigma=\mathbf1\oplus\psi,
\]
\[
 F_\sigma^{\sigma\sigma\sigma}
 =\frac1{\sqrt2}
 \begin{pmatrix}1&1\\1&-1\end{pmatrix},
 \qquad
 R_{\mathbf1}^{\sigma\sigma}=e^{-\pi i/8},
 \qquad
 R_\psi^{\sigma\sigma}=e^{3\pi i/8}.
\]
Ces matrices satisfont les relations de cohérence dans les conventions
indiquées. Le dédoublement
\(\mathcal C\boxtimes\mathcal C^{\rm rev}\) conserve les deux orientations et
évite d'introduire une chiralité par un choix externe.

\subsection{Persistance finie des états virtuels et absence de masse scalaire universelle}

Un état virtuel est une section de persistance finie dans une amplitude ou une
résolvante. Il n'appartient pas au spectre asymptotique et n'exige pas de masse de pôle
réelle. Ses paramètres sont la virtualité, le canal et la durée de
persistance. Cette définition évite de transformer une ligne interne de calcul en
une nouvelle espèce.

\section{Gravitation, torsion et critère spectral du mode collectif de spin 2}

\subsection{Exclusion tachyonique par positivité et séparation de la gravitation primaire par connexion et torsion}

La positivité exclut une espèce élémentaire avec \(m^2<0\) dans le
codomaine massique. Une valeur propre négative du hessien d'un fond ne définit pas un
tachyon comme particule : elle diagnostique l'instabilité du fond et exige de changer
la solution autour de laquelle on linéarise.

La gravitation entre dans la mécanique dimensionnelle par la connexion, la courbure et la torsion. Elle ne requiert pas de particule élémentaire de spin \(2\) comme donnée de départ.

\subsection{Pôle physique du hessien transverse sans trace et positivité du résidu}

Soit \(\mathcal S_{\rm grav}\) l'action géométrique et \(\mathcal H_{\rm TT}\) le sous-espace transverse et sans trace de son hessien en une solution. L'observable linéarisé est
\[
 \mathcal K_{\rm TT}
 =P_{\rm TT}\,\delta^2\mathcal S_{\rm grav}\,P_{\rm TT}.
\]
Un mode collectif de spin \(2\) n'existe comme particule asymptotique que
si la résolvante de \(\mathcal K_{\rm TT}\) possède un pôle physique de résidu
positif et respecte les contraintes. Si un tel pôle n'existe pas, l'interaction
gravitationnelle demeure intégralement dans la connexion et la torsion. Ainsi, le
graviton n'est pas un élément nécessaire de la loi des masses, et son apparition éventuelle
est décidée par un critère spectral, et non par une case nominale.

\section{Équivalence entre section persistante, limite inverse et représentation finie de particule}

\subsection{Quotient par équivalence de queue et compatibilité des familles de sections}

La définition généalogique de la particule et sa représentation opératorielle sont
deux présentations du même objet. Soit
\[
 \mathfrak r_K=(x_K,\gamma_K,B_K,s_K,n_K,U_K)
\]
un représentant fini à la profondeur \(K\), avec état, route, frontière,
orientation, compteur et représentation, et soient
\(r_{L,K}:\mathfrak R_L\to\mathfrak R_K\) les applications de raffinement. Une
section persistante est une famille
\[
 X=(\mathfrak r_K)_{K\ge K_0},
 \qquad r_{L,K}(\mathfrak r_L)=\mathfrak r_K
 \quad(L\ge K\ge K_0).
\]
Deux familles sont équivalentes si elles coïncident sur une queue cofinale.

\subsection{Condition de prolongement unique pour reconstruire une section à partir d'une évaluation finie}

\begin{theorem}[Équivalence de la section persistante et de la limite projective]
La catégorie des sections persistantes, définies comme des familles compatibles,
est canoniquement équivalente à
\[
 \varprojlim_K\mathfrak R_K/\!\sim_{\rm cola}.
\]
La projection à une profondeur finie est une évaluation, et non un foncteur inverse.
Un représentant \(\mathfrak r_K\) détermine une unique section persistante si
et seulement s'il possède un unique prolongement compatible modulo l'équivalence de queue ;
cette unicité est obtenue, par exemple, lorsque les applications de raffinement sont
bijectives sur la queue survivante à partir de \(K\).
\end{theorem}

\begin{proof}
Par définition, une section persistante est précisément un point de la limite
projective ; le quotient identifie les familles qui coïncident sur une queue cofinale.
L'évaluation est la projection canonique sur une composante. Sa fibre contient
tous les prolongements compatibles du représentant évalué, de sorte qu'elle
est unitaire exactement sous la condition de prolongement unique indiquée.
\end{proof}

Une particule HMT--MD est donc une classe de sections persistantes sans
obstruction terminale finie, munie d'une représentation, d'une orientation, d'une charge,
d'un spin et d'un caractère spectral compatibles avec le raffinement. La route est son
histoire observable, et non son identité réduite à une cellule. L'antiparticule
est la section conjuguée par l'involution dodécaphasique et la représentation
conjuguée. Une section finie obstruée est virtuelle ; un produit cocyclique de
sections est composé ; une représentation tressée persistante est anyonique.

\section{Théorème de clôture de la réalisation équivariante et classification exhaustive des espèces matérielles}

\subsection{Atlas exhaustif de 324 voies : 49 invariants par section et conservation du secteur nul}

\begin{theorem}[Clôture par codomaine physique]
Supposons construits l'atlas \(81/56/13/324\), le secteur nul, l'échelle
action--horloge--masse, les lecteurs \(K_D,K_\Omega\), le descripteur physique et la
représentation interne de chaque état. Alors tout descripteur admissible et
réalisable, c'est-à-dire tout \(X\) avec
\(d_{\mathcal P}(X)\in\operatorname{im}d_{\mathcal R}\),
détermine, sans consulter sa valeur externe :
\begin{enumerate}
\item une fibre, une orbite ou une multisection \(\mathscr S(X)\) ;
\item la paire positive bidirectionnelle et sa PVM conjointe
      \((\widehat M_X^{\beta,D},\widehat M_X^{\beta,\Omega},
      E_X^{D,\Omega})\) ;
\item une décomposition en blocs irréductibles ;
\item une mesure spectrale pour chaque couplage physique ;
\item un scalaire dans un secteur comprimé de rang un ou un multiplet dans un
      secteur réductible ; dans un secteur ouvert, un opérateur effectif et, seulement s'il
      existe un prolongement méromorphe avec un zéro isolé admissible, un pôle
      complexe ;
\item un objet catégorique lorsque la réalisation est topologique et ne possède pas encore
      de gap hamiltonien ;
\item un critère d'inexistence lorsque le type demandé contredit
      la positivité, la jauge ou l'asymptoticité.
\end{enumerate}
Aucun de ces résultats ne requiert de sélectionner une route au moyen de la masse
qui sera ensuite confrontée.
\end{theorem}

\begin{proof}
Le produit fibré construit la fibre de manière prospective. Le caractère et les
deux lecteurs construisent la paire positive et sa PVM conjointe. Lorsque le type
physique déclare une fonctionnelle scalaire, le calcul fonctionnel produit l'opérateur
correspondant ; l'espérance conditionnelle le projette sur le commutant de la
symétrie. La décomposition isotypique produit des scalaires ou des blocs par le lemme
de Schur. Les projecteurs de préparation produisent des mesures spectrales. La
réduction de Feshbach produit l'opérateur effectif ; le prolongement méromorphe
et un zéro isolé admissible produisent un pôle, s'il existe. Les secteurs
topologiques conservent leurs données de fusion et de tressage jusqu'à ce qu'il existe une
réalisation hamiltonienne. La positivité et les contraintes de jauge
excluent les types incompatibles. Chaque étape dépend uniquement de données internes ou
du couplage déclaré, jamais de la colonne externe.
\end{proof}

\paragraph{Contrôles de réfutation.}

La construction est réfutée dans son domaine si l'un quelconque des
faits suivants se produit :

\begin{enumerate}
\item la modification des valeurs externes change la fibre \(\mathscr S(X)\) ;
\item l'attribution cesse d'être équivariante sous l'inversion cellulaire,
      la conjugaison ou la jauge ;
\item une cellule ponctuelle est sélectionnée dans une multisection non centrale sans
      couplage brisant la symétrie ;
\item la synthèse de \(K_D\) et \(K_\Omega\) privilégie une orientation ou ne
      retrouve pas la diagonale lorsque les deux lecteurs coïncident ;
\item une compression à plusieurs valeurs propres est remplacée par une moyenne sans
      projecteur physique ;
\item le quotient de couleur produit des masses différentes pour le rouge, le vert et le bleu ;
\item des masses de quarks sont comparées sans déclarer le même schéma et la même échelle ;
\item une phase de rang un est présentée comme mélange PMNS après avoir échoué à la
      confrontation angulaire de rang complet ;
\item la neutralité est utilisée à elle seule pour affirmer le caractère Majorana ;
\item la masse d'un composé est obtenue en additionnant les masses des constituants sans
      registre de liaison ;
\item la largeur est insérée comme coordonnée de la signature au lieu d'être extraite
      d'un mode ouvert ou d'un pôle ;
\item une valeur propre négative du Hessien est publiée comme particule tachyonique ;
\item la gravitation est rendue dépendante d'un graviton élémentaire introduit
      avant la connexion, la courbure et la torsion.
\end{enumerate}

\subsection{Inventaires complets de 471 masses et 384 largeurs et reconnaissance métrologique ultérieure}

La comparaison externe n'est ouverte qu'après avoir figé le descripteur,
la fibre, l'opérateur, le projecteur, le schéma et l'échelle. Chaque ligne doit distinguer la masse
stable, la masse de pôle, l'estimateur de résonance et la largeur. L'atlas de \(324\)
fiches et les inventaires de \(471\) masses et de \(384\) largeurs constituent
l'exposé exhaustif qui suit. La construction conceptuelle et chaque fiche
renvoient réciproquement l'une à l'autre par leurs identifiants de route. Les fiches
conservent tous leurs champs et sont composées comme des unités cohérentes ; elles ne sont pas
remplacées par des tableaux historiques abrégés.

La comparaison ne consiste pas à élargir une liste de noms, mais à appliquer,
pour chaque type physique, l'application complète qui conduit de la généalogie de route
à son codomaine observable.
